\documentclass[12pt,a4paper]{article}
% ============================================================
%  Structural Persistence Series — Shared Preamble
% ============================================================
\usepackage{fontspec}
\setmainfont{Hiragino Mincho ProN}
\setsansfont{Hiragino Sans}
\setmonofont{Menlo}

\usepackage{iftex}
\ifXeTeX
  \XeTeXlinebreaklocale "ja"
  \XeTeXlinebreakskip=0pt plus 1pt minus 0.1pt
\fi

\usepackage[margin=22mm]{geometry}
\usepackage{amsmath,amssymb,amsthm}
\usepackage{xcolor}
\usepackage{graphicx}
\usepackage{hyperref}
\usepackage{booktabs}
\usepackage{fancyhdr}
% enumitem not available in this TeX installation

% --- Colours ------------------------------------------------
\definecolor{PaperAccent}{HTML}{1F4E79}
\definecolor{PaperGray}{HTML}{51606F}

\hypersetup{
  colorlinks=true,
  linkcolor=PaperAccent,
  urlcolor=PaperAccent,
  citecolor=PaperAccent,
  pdflang=ja
}
\urlstyle{same}

% --- Typography ---------------------------------------------
\setlength{\parindent}{1.5em}
\setlength{\parskip}{0pt}
\setlength{\emergencystretch}{3em}
\linespread{1.08}
\setlength{\abovedisplayskip}{8pt plus 2pt minus 4pt}
\setlength{\belowdisplayskip}{8pt plus 2pt minus 4pt}
\setlength{\abovedisplayshortskip}{6pt plus 2pt minus 2pt}
\setlength{\belowdisplayshortskip}{6pt plus 2pt minus 2pt}

% --- Section label names ------------------------------------
\renewcommand{\abstractname}{要旨}

% --- Theorem-like environments ------------------------------
\theoremstyle{plain}
\newtheorem{proposition}{命題}[section]
\newtheorem{lemma}[proposition]{補題}
\newtheorem{corollary}[proposition]{系}

\theoremstyle{definition}
\newtheorem{definition}{定義}[section]
\newtheorem{assumption}{条件}[section]

\theoremstyle{remark}
\newtheorem*{remark}{注}

% --- Running header -----------------------------------------
\pagestyle{fancy}
\fancyhf{}
\renewcommand{\headrulewidth}{0.4pt}
\fancyhead[L]{\small\textcolor{PaperGray}{\nouppercase{\leftmark}}}
\fancyhead[R]{\small\textcolor{PaperGray}{\thepage}}
\fancypagestyle{plain}{%
  \fancyhf{}
  \renewcommand{\headrulewidth}{0pt}
  \fancyfoot[C]{\small\textcolor{PaperGray}{\thepage}}
}

% --- Title block macro --------------------------------------
\newcommand{\PaperTitleBlock}[3]{%
  % #1 = title, #2 = subtitle, #3 = date
  \thispagestyle{plain}%
  \begin{center}
  {\LARGE\bfseries #1\par}
  \vspace{0.4em}
  {\normalsize #2\par}
  \vspace{1.0em}
  {\large Akihito Sunagawa\par}
  \vspace{0.3em}
  {\normalsize #3\par}
  \end{center}
  \vspace{1.6em}
}

% Legacy 2-arg version (backward compat)
\newcommand{\PaperTitleBlockSimple}[2]{%
  \PaperTitleBlock{#1}{}{#2}
}


\begin{document}

\PaperTitleBlock
  {構造持続の収支原理}
  {— 消耗と回復の差が構造維持可能性を支配する —}
  {2026年4月12日}

\begin{abstract}
構造持続の最小形式は、回復を明示しない収縮モードとして

\[
S = M e^{-L}
\]

を与える。ここで \(L\) は、構造維持可能な状態集合がどれだけ縮小したかを対数比で測った累積構造消耗量である。

本稿は、この最小核を、回復・修復・冗長化・外部支援・ロールバックを含む形式へ拡張する。各時点の構造消耗量を \(d_{t}\)、回復量を \(r_{t}\) とし、その差

\[
b_{t} := d_{t}-r_{t}
\]

を純消耗量と呼ぶ。有限時間地平 \(n\) までの累積純消耗量を

\[
B_{n} := \sum_{t<n} b_{t}
\]

と置けば、構造維持可能領域の質量は

\[
m(V^{(n)}) = m(V^{(0)}) e^{-B_{n}}
\]

と書ける。有効維持余力を明示する読者向けの形式では

\[
S_{n} = M_{n} e^{-B_{n}}
\]

である。

この式は、構造持続の最小核 \(S=Me^{-L}\) を置き換えるものではない。回復量がない場合、すなわち \(r_{t}=0\) の場合には \(B_{n}=L_{n}\) となり、最小形式に戻る。本稿の中心主張は、回復を含む系では、構造持続を決めるのは消耗量そのものではなく、消耗量と回復量の差である、という一点にある。

本稿の意義は、回復を含む系でも指数核は失われず、消耗量ではなく累積純消耗量 \(B\) が構造維持可能性を支配することを示す点にある。
\end{abstract}

\section{問い}\label{ux554fux3044}

構造は削られるだけではない。ソフトウェアには rollback や冗長化があり、組織には手順、引き継ぎ、権限分離がある。長期対話には範囲づけや外部整理があり、継続学習には replay、adapter 分離、依存関係に沿った再同期がある。

したがって、回復を含む系で問うべき量は、「どれだけ削られたか」だけではない。

構造維持可能性はどれだけ消耗し、それに対してどれだけ回復したか。

本稿では、この差し引きを構造持続の収支原理と呼ぶ。ここでの収支は均衡を意味しない。構造消耗と回復を同じ尺度で差し引く accounting を意味する。消耗が回復を上回れば構造は崩壊、転移、停止へ向かいやすくなり、回復が消耗を上回れば維持または回復へ向かう。

したがって、本稿でいう消耗、維持、回復の局面は、系が均衡状態にあるかどうかの判定ではない。後に述べるように、それらは事前固定された観測単位の上で、差し引き量の符号を読むための局所的な会計分類である。

\section{二段階更新}\label{ux4e8cux6bb5ux968eux66f4ux65b0}

系 \(X\) と、各時刻 \(t\) における構造維持可能集合

\[
V^{(t)} \subseteq X
\]

を考える。\(m\) は、各 \(V^{(t)}\) を比較するために事前固定された有限測度である。対象となる構造条件、測度、観測単位、時間地平は、最小形式と同じく事前に固定されているものとする。

各時刻の更新を、観測単位の内部で二段階に分ける。

\[
\begin{aligned}
V_{t}^- :&= K_{t}(V^{(t)}), \\
\qquad \\
V^{(t+1)} :&= R_{t}(V_{t}^-).
\end{aligned}
\]

ここで \(K_{t}\) は構造維持可能領域を削る作用、\(R_{t}\) はその領域を再拡大する作用である。現実の系で消耗と回復が同時に起きる場合でも、観測単位を固定すれば、その内部をこの二段階更新として読むことができる。

以下では、対数比を取る各段階で \(m(V^{(t)})\), \(m(V_{t}^-)\), \(m(V^{(t+1)})\) は正かつ有限であるとする。質量がゼロに到達する場合は、有限の対数比更新ではなく、構造維持可能領域が尽きる collapse boundary として別に扱う。

\section{消耗量・回復量・純消耗量}\label{ux6d88ux8017ux91cfux56deux5fa9ux91cfux7d14ux6d88ux8017ux91cf}

各時刻の構造消耗量を

\[
d_{t} := -\log \frac{m(V_{t}^-)}{m(V^{(t)})}
\]

と定める。各時刻の回復量を

\[
r_{t} := \log \frac{m(V^{(t+1)})}{m(V_{t}^-)}
\]

と定める。どちらも同じ対数尺度上の量である。

純消耗量を

\[
b_{t} := d_{t}-r_{t}
\]

と定義する。このとき中間集合 \(V_{t}^-\) の質量は打ち消し合い、

\[
\begin{aligned}
b_{t} \\
&= \\
-\log \frac{m(V^{(t+1)})}{m(V^{(t)})}
\end{aligned}
\]

が成り立つ。

すなわち、各時刻の局所的な変化は、消耗と回復を別々に導入しても、最終的には始点と終点の対数比として表される。

\section{累積収支と指数核}\label{ux7d2fux7a4dux53ceux652fux3068ux6307ux6570ux6838}

累積純消耗量を

\[
B_{n} := \sum_{t=0}^{n-1} b_{t}
\]

と置く。このとき、望遠鏡積により

\[
m(V^{(n)}) = m(V^{(0)}) e^{-B_{n}}
\]

が成り立つ。

証明。 前節の式より

\[
\begin{aligned}
e^{-b_{t}} \\
&= \\
\frac{m(V^{(t+1)})}{m(V^{(t)})}.
\end{aligned}
\]

したがって

\[
\begin{aligned}
e^{-B_{n}} \\
&= \\
\prod_{0\le t<n} e^{-b_{t}} \\
&= \\
\prod_{0\le t<n} \\
\frac{m(V^{(t+1)})}{m(V^{(t)})} \\
&= \\
\frac{m(V^{(n)})}{m(V^{(0)})}.
\end{aligned}
\]

よって

\[
m(V^{(n)}) = m(V^{(0)})e^{-B_{n}}.
\]

証明終。

資源項を明示する読者向けの形式では、

\[
S_{n} = M_{n} e^{-B_{n}}
\]

と書く。対象期間が文脈上固定されているときは

\[
S = M e^{-B}
\]

と略記する。

ここで \(S_{n}\) は、対象となる構造条件に対する構造持続ポテンシャルである。\(S_{n}=0\)、または事前固定された閾値 \(S_{n}\le S_{c}\) への到達は、その構造条件に関する構造維持境界への到達を意味する。その観測上の現れは、ドメインに応じて崩壊、機能不全、停止、相転移、構造再編として現れうる。

\(M_{n}\) は、上の経路ごとの指数核から導かれる量ではない。構造維持に利用できる有効維持余力を表すリソース側のスカラーである。構造維持境界への到達には、累積純消耗量 \(B_{n}\) が表す維持可能領域側の経路、\(M_{n}=0\) が表すリソース側の経路、またはその両方がある。本稿は \(M_{n}\) の成分分解や普遍的な予測力を主張しない。

\section{最小形式との関係}\label{ux6700ux5c0fux5f62ux5f0fux3068ux306eux95a2ux4fc2}

構造持続の最小形式では、回復量を明示しない。すなわち、各時刻で

\[
r_{t}=0
\]

と置く。このとき

\[
\begin{aligned}
b_{t}&=d_{t}, \\
\qquad \\
B_{n}&=\sum_{t<n}d_{t}=L_{n}
\end{aligned}
\]

である。

したがって

\[
S_{n}=M_{ne}^{-B_{n}}
\]

は

\[
S_{n}=M_{ne}^{-L_{n}}
\]

に戻る。

つまり、最小形式 \(S=Me^{-L}\) は、収支原理 \(S=Me^{-B}\) の特例である。逆に言えば、収支原理は、最小形式の指数核を保ったまま、回復量を差し引くための拡張である。

回復を入れても指数形は変わらない。変わるのは指数の中身である。回復を明示しない場合には累積構造消耗量 \(L\) が入り、回復を含む場合には累積純消耗量 \(B\) が入る。

\section{三つの局面}\label{ux4e09ux3064ux306eux5c40ux9762}

収支原理では、崩壊、維持、回復は別々の現象ではなく、同じ符号付き量 \(b_{t}\) の三つの読みとして現れる。

\[
b_{t}>0
\]

なら消耗局面である。これが続けば \(B_{n}\) は増え、構造維持可能領域は縮小する。

\[
b_{t}=0
\]

なら維持局面である。ただし、ここでいう維持は局所的な釣り合いであって、理論全体を均衡理論へ還元するものではない。

\[
b_{t}<0
\]

なら回復局面である。これが続けば \(B_{n}\) は減り、構造維持可能性は回復方向へ向かう。

§1 で述べた通り、この局面分類は均衡判定ではなく、事前固定された観測単位の上での会計分類である。隣接する二つの区間を併合すれば

\[
\begin{aligned}
b_{t:t+2} \\
&= \\
-\log \frac{m(V^{(t+2)})}{m(V^{(t)})} \\
&= \\
b_{t}+b_{t+1}
\end{aligned}
\]

となり、累積純消耗量 \(B_{n}\) は加法的に整合する。しかし、細かい粒度では消耗局面の直後に回復局面が現れる経路が、粗い粒度では維持局面として読まれることがある。したがって、消耗・維持・回復という局面分類や hitting boundary を経験的に主張する場合には、観測単位、段階列、時間地平を事前に固定しなければならない。

確率過程として扱う場合には、\(\mathbb E[b_{t}]\) の符号が期待値レベルの傾向を与える。ただし、期待値レベルの傾向から高確率の崩壊境界が無条件に従うわけではない。高確率境界、停止時刻、hitting-time bound を得るには、bounded increments、MGF、Chernoff / KL、margin 条件などの追加構造が必要である。

\section{適用と限界}\label{ux9069ux7528ux3068ux9650ux754c}

本稿の核は、事前固定された構造維持問題に対する経路ごとの代数核である。すなわち、対象となる構造条件、測度、時間地平、観測単位、消耗量、回復量を固定したとき、

\[
m(V^{(n)}) = m(V^{(0)}) e^{-B_{n}}
\]

が成り立つ、という主張である。

現実ドメインでは、何を維持対象とし、どの測度で比較し、何を \(d_{t}\) や \(r_{t}\) の指標にするかを探索的に発見する段階がある。この探索は候補生成であり、その同じデータ上の説明は support ではない。support は、写像、指標、baseline、metric、split、判定規則を凍結した後の holdout、future surface、fresh archive、または outside rerun でのみ成立する。

また、本稿は次を主張しない。第一に、あらゆるドメインで自然な \(V,m,d_{t},r_{t}\) が一意に定まるとは主張しない。第二に、観測指標が真の \(r_{t}\) を直接測定しているとは主張しない。第三に、期待値レベルの傾向から高確率境界や因果主張が無条件に従うとは主張しない。第四に、構造持続理論がこの式だけで普遍法則として確立したとは主張しない。

これらの限界は、理論を弱めるためではなく、空虚化を防ぐための境界である。詳細な写像手順、support 判定、観測可能性の層分類、既存理論との接続は補論に置く。本稿の役割は、最小核 \(S=Me^{-L}\) が、回復を含む \(S=Me^{-B}\) へどう持ち上がるかを短く閉じることである。

\section{結論}\label{ux7d50ux8ad6}

構造持続の最小形式は

\[
S = M e^{-L}
\]

を与える。本稿の収支原理は、それを

\[
S = M e^{-B}
\]

へ拡張する。ここで \(B\) は、消耗量から回復量を差し引いた累積純消耗量である。

この置き換えにより、崩壊、維持、回復を同じ座標で読むことができる。回復がなければ \(B=L\) であり、最小形式に戻る。回復があれば、構造持続は単なる消耗の蓄積ではなく、消耗と回復の収支として決まる。

この同じ座標は、ドメイン横断の写像と設計転用の基礎にもなる。あるドメインで見つかった消耗の局所化、回復経路の保存、余力の確保は、別ドメインの検証可能な介入候補へ翻訳できる。ただし、その転用は support の輸入ではなく、補論で定める写像手順と凍結検証を要する。

本稿の中心は、この一点である。構造は、存在しているかどうかだけで決まるのではない。維持できる余地がどれだけ消耗し、どれだけ回復され、どれだけ純消耗として残るかで決まる。このため収支原理は、最小形式を保ったまま、閉じた縮小の理論を回復を含む構造維持の理論へ拡張する。
\end{document}
